\documentclass[11pt]{article}

%------------------------------------------------
% Packages
%------------------------------------------------
\usepackage[a4paper,margin=2.5cm]{geometry}
\usepackage{amsmath,amssymb,amsthm}
\usepackage{mathtools}
\usepackage{enumitem}
\usepackage{hyperref}

\hypersetup{
    colorlinks=true,
    linkcolor=blue,
    citecolor=blue,
    urlcolor=blue
}

%------------------------------------------------
% Title Information
%------------------------------------------------

\title{\bfseries How to make (Bregman-)Kaczmarz methods fast}

\author{
Dirk Lorenz\thanks{Center for Industrial Mathematics, University of Bremen,
\texttt{d.lorenz@uni-bremen.de}}
}

\date{}

%------------------------------------------------

\begin{document}

\maketitle

\vspace{-1em}

\noindent
\textbf{Keywords:}
constrained optimization, Kaczmarz methods

\vspace{0.5em}


\vspace{1em}

%%%%%%%%%%%%%%%%%%%%%%%%%%%%%%%%%%%%%%%%%%%%%%%%%%
\section*{Abstract}

Bregman-Kaczmarz methods can compute solutions of systems of linear equations that additionally minimize a strongly convex functions (e.g. minimum-norm solutions, but also sparse solutions). Their plain version is oftentimes quite slow and hence, there is interest in speeding up the methods. Optimal step-sizes are, on the one hand, often hard to compute and, on the other hand, do not always bring much improvements for. In this talk I will treat several acceleration mechanisms e.g. based on momentum (with a computable approximate minimal error criterion) and dual acceleration.

%%%%%%%%%%%%%%%%%%%%%%%%%%%%%%%%%%%%%%%%%%%%%%%%%%
\vspace{1em}

\noindent
\textbf{Acknowledgements.}
This work was supported by the International Training Network TraDE-OPT fundedd, MSCA Grant agreement ID: 861137

%%%%%%%%%%%%%%%%%%%%%%%%%%%%%%%%%%%%%%%%%%%%%%%%%%
\vspace{1em}

\begin{thebibliography}{9}

\bibitem{1}
D.A. Lorenz, M. Winkler.
\newblock  Minimal error momentum Bregman-Kaczmarz
\newblock {\em Linear Algebra and its Applications} 709, 416-448, 2025

\bibitem{2}
L. Tondji, I. Necoara, D.A. Lorenz.
\newblock Acceleration and restart for the randomized Bregman-Kaczmarz method
\newblock {\em Linear Algebra and its Applications} 699, 508-538, 2024

\bibitem{3}
L. Tondji, D.A. Lorenz.
\newblock Faster randomized block sparse Kaczmarz by averaging
\newblock {\em Numerical Algorithms} 93 (4), 1417-1451, 2023

\end{thebibliography}

\end{document}
